\chapter{Fourier and Mellin Comparisons in the Explicit Formula}
\label{v4-arith:w7-i:chapter}

The Guinand and Weil expressions use the same gamma density with
opposite placement in the explicit formula. Fourier inversion gives
a subtracted integral in the logarithmic variable. Poisson summation
uses the Fourier transform of its test, whereas multiplicative
inversion is a different operation.

\section{The signed convention dictionary}
\begin{theorem}[the compact Guinand formula]
\label{v4-arith:w7-i:thm:guinand-explicit}
For $f\in\mathscr A$ put $h=\mathcal F_+f$. Then
\[
 \sum_\rho h(z_\rho)=h(i/2)+h(-i/2)
       +\int_{\mathbb R}h(u)\Phi_\infty^G(u)du
       -\sum_{n\ge2}\frac{\Lambda(n)}{\sqrt n}
          \bigl(f(\log n)+f(-\log n)\bigr),
\]
where
\[
 \Phi_\infty^G(u)=\frac{\Re\psi_\Gamma(1/4+iu/2)-\log\pi}{2\pi}.
\]
The zero sum is absolutely convergent and the prime sum is finite.
\end{theorem}
\begin{proof}
This is Theorem~\ref{v4-arith:pm:weil}, with
$z_\rho=-i(\rho-1/2)$. In particular, the zero argument retains
its imaginary part when $\Re\rho\ne1/2$.
\end{proof}

\begin{proposition}[exact kernel dictionary]
\label{v4-arith:w7-i:prop:weil-vs-guinand}
If $\Phi_\infty^W=k/(2\pi)$ denotes the density of $W_\infty$, then
\[
 \Phi_\infty^G=-\Phi_\infty^W.
\]
Their difference is $-k/\pi$, which is nonconstant.
For the auxiliary density $g/(2\pi)$, the relation is instead
$\Phi_\infty^G-g/(2\pi)=-\log\pi/\pi$.
\end{proposition}
\begin{proof}
Insert $k=\log\pi-\Re\psi_\Gamma$ and
$g=\log\pi+\Re\psi_\Gamma$. Strict radial monotonicity of $k$
proves nonconstancy. The auxiliary density is not the density of
$W_\infty$.
\end{proof}

\section{The exact time-variable gamma integral}
\begin{theorem}[subtracted gamma pairing]
\label{v4-arith:w7-i:prop:dual-kernel}
For every $f\in\mathscr A$,
\[
 \begin{split}
 W_\infty(f)={}&(\log\pi+\gamma)f(0)\\
 &+\int_0^\infty
 \frac{e^{-t/2}(f(t)+f(-t))-2e^{-2t}f(0)}{1-e^{-2t}}dt.
 \end{split}
\]
The integral is absolutely convergent, including at zero.
For tests vanishing near zero it reduces to the pairing against
$e^{-|t|/2}/(1-e^{-2|t|})$ away from the origin.
\end{theorem}
\begin{proof}
Use the integral
\[
 \psi_\Gamma(z)+\gamma
 =\int_0^\infty\frac{e^{-v}-e^{-zv}}{1-e^{-v}}dv,
 \qquad\Re z>0,
\]
from \href{https://dlmf.nist.gov/5.9.E16}{DLMF 5.9.16}.
Truncate first to $\delta\le v\le M$. Fourier inversion gives
\[
 \frac1{2\pi}\int h(u)\cos(uv/2)du
 =\frac{f(v/2)+f(-v/2)}2.
\]
The numerator in the digamma integrand is bounded near zero by
$Cv(1+|u|)$; the Schwartz decay of $h$ gives an integrable bound
after division by $1-e^{-v}$. Exponential decay handles infinity.
Dominated convergence therefore removes both cutoffs.
Set $v=2t$. The resulting numerator is $3tf(0)+O(t^2)$ near zero,
while its denominator is $2t+O(t^2)$. At infinity the compact tests
vanish and the subtracted term decays exponentially.
\end{proof}

\section{Poisson reciprocity and multiplicative inversion}
For this section only, put
$\widehat\varphi_{2\pi}(u)=\int\varphi(x)e^{-2\pi iux}dx$.
This differs from the arithmetic transform by
$\widehat\varphi_{2\pi}(u)=\mathcal F_+\varphi(-2\pi u)$.

\begin{proposition}[the zero-moment Poisson identity]
\label{v4-arith:w7-i:thm:burnol-copoisson}
Let $\varphi$ be an even Schwartz function with $\varphi(0)=0$ and
$\int_0^\infty\varphi(x)dx=0$.
Set $C\varphi(y)=\sum_{n\ge1}\varphi(ny)$ for $y>0$.
Then
\[
 C\varphi(y)=y^{-1}C\widehat\varphi_{2\pi}(1/y).
\]
Both sides decay faster than every power at zero and infinity.
The Mellin integral of $C\varphi$ is entire, and on $\Re s>1$,
\[
 \int_0^\infty C\varphi(y)y^{s-1}dy
 =\zeta(s)\int_0^\infty\varphi(y)y^{s-1}dy.
\]
\end{proposition}
\begin{proof}
The two zero moments imply
$\widehat\varphi_{2\pi}(0)=0$ and
$\int_0^\infty\widehat\varphi_{2\pi}(u)du=\varphi(0)/2=0$.
Poisson summation, with its zero terms removed and even terms paired,
gives the first identity. Schwartz estimates give rapid decay at
infinity. The same estimates for the Fourier transform and the
identity give rapid decay at zero. The Mellin integral therefore
converges locally uniformly for all $s$.
For $\Re s>1$, absolute convergence permits exchanging the sum and
integral, and substitution $x=ny$ gives the Dirichlet series $\zeta(s)$.
\end{proof}

\begin{proposition}[reciprocal substitution is not Poisson reciprocity]
Replacing $\widehat\varphi_{2\pi}(t)$ in the preceding identity by
$t^{-1}\varphi(1/t)$ gives a false assertion, even for the stated
zero-moment tests.
\end{proposition}
\begin{proof}
Choose an even compact smooth $\varphi$ supported on the positive
side in $(2.2,2.6)$, with integral zero and $\varphi(2.4)=1$.
Two disjoint bumps in that interval, with adjusted integrals, give
such a function. Then $C\varphi(2.4)=1$.
With $\varphi^\vee(t)=t^{-1}\varphi(1/t)$ for $t>0$, the proposed
right side equals
\[
 \frac1{2.4}\sum_{n\ge1}\varphi^\vee(n/2.4)
 =\sum_{n\ge1}\frac1n\varphi(2.4/n)=1.
\]
At $y=1.2$, however, $C\varphi(1.2)=\varphi(2.4)=1$, while
$\sum_{n\ge1}n^{-1}\varphi(1.2/n)=0$.
This contradicts the proposed identity at $y=1.2$.
\end{proof}

\section{The arithmetic comparison}
The signed zero map $E$ and the compressed translation operator
$B_R$ give the exact comparison
\[
 \langle Ef,JEf\rangle
 =P(f*f^*)-W_\infty(f*f^*)-\langle B_Rf,f\rangle.
\]
Poisson reciprocity does not remove either polar moment or change
the sign of the gamma term. A proposed positive comparison must
preserve this identity on its stated test space.
